\documentclass{article}

\usepackage[a4paper, margin=1in]{geometry}
\usepackage[french]{isodate}
\usepackage{setspace}

\title{Comptes-rendus 2}
\author{Aden Perry}
\date{\today}

\doublespacing

\begin{document}
\maketitle

Anjali Prakash-Ream --- Les exoplanètes

\begin{enumerate}
	\item \textbf{Des exoplanètes gazeux:} Les exoplanètes gazeux sont très
	      courant, mais j'étais hereux d'être rappelé.
	\item \textbf{Des exoplanètes verglacé:} Les exoplanètes sont aussi très
	      courant. Les exoplanètes verglacé sont ma favorite, alors j'étais
          excité de les voir.
	\item \textbf{Des exoplanètes rocheux:} Les exoplanètes rocheux sont une
	      variante moins courant, mais le plus similaire à la Terre. J'etais
	      hereux d'apprendre de cets exoplanètes parce que j'ai oublié le
	      classifcation.
	\item \textbf{La methode de détection de vitesse radiale:} Les exoplanètes
	      peuvent être détecté par utiliser les diffèrences d'émission
	      photonique et l'effet Doppler. Je pense qu'il est très cool que nous
	      peuvons voir les diffèrences, et en particulier parce que la methode
	      est très précis.
	\item \textbf{La methode de détection de transit:} Les exoplanètes
	      peuvent aussi être détecté par utiliser de transit. C'est une methode
	      aussi très cool, et je suis étonné que nous avons decouvert beaucoup
	      d'exoplanètes avec cette methode.
	\item \textbf{La methode de détection de imagerie directe:} Il n'est pas
	      une methode très flable, mais il est interessant parce que il nous
	      donné un image de l'exoplanète.
\end{enumerate}

\end{document}
